%% --------------------------------------------------------
\section{Інфрачервоні (ІЧ) спектри}
%% --------------------------------------------------------

\textbf{ІЧ-спектроскопія} досліджує взаємодію молекул з інфрачервоним випромінюванням.
Поглинання відбувається тоді, коли частота випромінювання збігається з частотою
внутрішнього коливання молекули, тобто виконується умова резонансу:
\[
    h\nu = \Delta E_{\text{колив.}}
\]

Ці переходи відповідають змінам \emph{коливальних} рівнів енергії.
Типові частоти для молекулярних коливань лежать у межах
\[
    400 - 4000~\text{см}^{-1},
\]
що відповідає довжинам хвиль у мікрометровому діапазоні ($2.5 - 25~\mu$m).

ІЧ-спектр відображає лише ті \textit{коливальні моди молекули}, які супроводжуються \emph{зміною її дипольного моменту}.

%% --------------------------------------------------------
\subsection{Інтенсивності ІЧ-поглинання}
%% --------------------------------------------------------

У спектроскопії інтенсивність смуги --- це \textbf{інтегральна поглинальна здатність}, тобто
площа під піком у спектрі.
Її фізичний зміст --- це міра того, \emph{наскільки сильно молекула поглинає ІЧ-випромінювання}
при переході $v=0 \to 1$.

Із теорії взаємодії електромагнітного поля з речовиною відомо, що інтенсивність $A_i$ для $i$-го коливального переходу визначається як:
\[
    A_i = \frac{N_A \pi}{3 c^2}
        \left| \frac{\partial \boldsymbol{\mu}}{\partial Q_i} \right|^2,
\]
де $\tfrac{\partial \boldsymbol{\mu}}{\partial Q_i}$ --- похідна дипольного моменту за нормальною координатою
$Q_i$ у рівноважному положенні, $N_A$ --- число Авогадро, а $c$ --- швидкість світла.

\paragraph{Фізичний зміст.}
Величина $A_i$ має розмірність площі на моль:
\[
    [A_i] = \text{см}^2 \text{/моль},
\]
тобто її можна трактувати як \textbf{ефективну площу поглинання на один моль молекул} речовини.
Чим більша ця площа --- тим інтенсивніше смуга в спектрі.

\paragraph{Зв’язок з експериментом.}
В експериментальній ІЧ-спектроскопії вимірюють \emph{оптичну густину} $A$ (безрозмірну величину),
яка пов’язана із концентрацією $c$ та довжиною кювети $l$ законом Бугера–Ламберта–Бера:
\[
    A = \varepsilon\, c\, l,
\]
де $\varepsilon$ --- молярний коефіцієнт поглинання.
Інтеграл від $\varepsilon(\tilde{\nu})$ за частотою $\tilde{\nu}$ визначає \textbf{молярну інтегральну інтенсивність смуги}:
\[
    I = \int \varepsilon(\tilde{\nu})\, d\tilde{\nu},
\]
яка має ті ж самі одиниці, що й $A_i$, тобто см$^2$/моль.

\paragraph{Одиниці км/моль.}
Щоб уникнути дуже малих чисел, цю площу традиційно виражають у \textbf{кілометрах на моль}:
\[
    1~\text{км/моль} = 10^5~\text{см}^2/\text{моль}.
\]
Тоді типові значення інтенсивностей для фундаментальних коливань молекул
лежать у межах 10–1000~км/моль, що зручно для порівняння та таблиць.

\paragraph{Числова форма.}
Після підстановки числових значень констант:
\[
    N_A = 6.022 \times 10^{23}~\text{моль}^{-1}, \quad
    c = 2.998 \times 10^{10}~\text{см/с},
\]
отримуємо:
\[
    \boxed{I = 974.8638 \times
    \left| \frac{\partial \boldsymbol{\mu}}{\partial Q} \right|^2_{\text{a.u.}} \text{ км/моль}.}
\]

\paragraph{Зв’язок між $A_i$, $I$ і коефіцієнтом $K$.}
Формально експериментальна інтегральна інтенсивність $I_i$
і теоретичне вираження $A_i$ описують одну й ту саму фізичну величину,
записану в різних одиницях:
\[
    I_i = K\,A_i.
\]
Коефіцієнт $K$ враховує перехід від атомних одиниць, у яких обчислюється
похідна $\tfrac{\partial \boldsymbol{\mu}}{\partial Q_i}$, до лабораторних одиниць (км/моль).
Підставивши фізичні сталі,
\[
    K = \frac{N_A \pi}{3 c^2} (e a_0)^2 \times 10^5,
\]
отримуємо числове значення
\[
    \boxed{K = 974.8638,}
\]
тому остаточний вираз для інтегральної інтенсивності має вигляд:
\[
    \boxed{I_i = 974.8638
    \left| \frac{\partial \boldsymbol{\mu}}{\partial Q_i} \right|^2_{\text{a.u.}}~
    [\text{км/моль}].}
\]



%% --------------------------------------------------------
\subsection{Класифікація за інтенсивністю}
%% --------------------------------------------------------

\begin{center}
\begin{tblr}{ll}
\toprule
\textbf{Інтенсивність (км/моль)} & \textbf{Класифікація} \\
\midrule
$I > 100$ & Сильна смуга \\
$10 < I \leq 100$ & Середня смуга \\
$I \leq 10$ & Слабка смуга \\
\bottomrule
\end{tblr}
\end{center}


%=========================================================
\section{Розрахунок ІЧ-інтенсивностей}
%=========================================================

ІЧ-інтенсивність нормальної моди пропорційна квадрату похідної дипольного моменту
за координатою цієї моди:
\[
I_i \sim \left| \frac{d\boldsymbol{\mu}}{dQ_i} \right|^2.
\]
На практиці цю похідну можна отримати двома різними шляхами, однак у

%---------------------------------------------------------
\subsection{Чисельне диференціювання}
%---------------------------------------------------------

У цьому підході координати атомів зсуваються вздовж $i$-ї нормальної моди
на малу величину~$\delta$, після чого обчислюються дипольні моменти
у двох геометріях --- \emph{«плюс»} і \emph{«мінус»}.
Похідна апроксимується центральною різницею:
\[
\left( \frac{d\boldsymbol{\mu}}{dQ_i} \right)_0 \approx
\frac{\boldsymbol{\mu}(Q_i+\delta) - \boldsymbol{\mu}(Q_i-\delta)}{2\delta}.
\]

В файлі \inputcode{h2o_ir_spectrum.py} наведено програму, яка  використовує \textbf{метод центральних різниць} для чисельного обчислення $\frac{d\boldsymbol{\mu}}{dQ_i}$:



\subsubsection{Алгоритм обчислення}

\begin{enumerate}
    \item \textbf{Отримання нормальної моди:}

    З діагоналізації масово-зваженого Гесіану отримуємо власний вектор $\mathbf{L}_k$ (нормальну моду):
    \begin{equation}
        \mathbf{L}_k = (l_{1x}, l_{1y}, l_{1z}, l_{2x}, l_{2y}, l_{2z}, \ldots, l_{Nx}, l_{Ny}, l_{Nz})^T
    \end{equation}
    де $N$ --- кількість атомів.

    \begin{minted}{python}
    freq_info = thermo.harmonic_analysis(mol, hess)
    normal_modes = freq_info['norm_mode']       # власні вектори

    L_k = normal_modes[:, k].reshape(natm, 3)
    \end{minted}

    Код \inlinecode{L\_k = normal\_modes[:, k].reshape(natm, 3)} відповідає $\mathbf{L}_k$ у формулі.

    \item \textbf{Зміщення координат:}

    Для чисельного диференціювання вибираємо крок $h$ та зміщуємо атоми вздовж нормальної моди:
    \begin{align}
        \mathbf{R}^{(+)} &= \mathbf{R}_0 + h \cdot \mathbf{L}_k \\
        \mathbf{R}^{(-)} &= \mathbf{R}_0 - h \cdot \mathbf{L}_k
    \end{align}

    \begin{minted}{python}
    h = 0.001  # Bohr
    R_0 = mol.atom_coords()
    R_pos = R_0 + h * L_k
    R_neg = R_0 - h * L_k
    \end{minted}

    \item \textbf{SCF-розрахунки:}

    Для кожної зміщеної геометрії виконуємо SCF-розрахунок, отримуємо енергію $E$ та дипольний момент $\boldsymbol{\mu}$:
    \begin{align}
        E^{(+)}, \boldsymbol{\mu}^{(+)} &= \text{SCF}(\mathbf{R}^{(+)}) \\
        E^{(-)}, \boldsymbol{\mu}^{(-)} &= \text{SCF}(\mathbf{R}^{(-)})
    \end{align}

    \begin{minted}{python}
    mf_pos = scf.RHF(mol_pos); mf_pos.kernel()
    dip_pos = mf_pos.dip_moment(unit='Bohr')

    mf_neg = scf.RHF(mol_neg); mf_neg.kernel()
    dip_neg = mf_neg.dip_moment(unit='Bohr')
    \end{minted}

    \item \textbf{Обчислення похідної дипольного моменту:}

    Чисельна похідна за нормальною координатою $Q_k$:
    \begin{equation}
        \frac{\partial \boldsymbol{\mu}}{\partial Q_k} = \frac{\boldsymbol{\mu}^{(+)} - \boldsymbol{\mu}^{(-)}}{2h}
    \end{equation}

    \begin{minted}{python}
    dip_derivative = (dip_pos - dip_neg) / (2 * h)
    \end{minted}

    \item \textbf{Обчислення інтегральної інтенсивності:}

    Інтенсивність ІЧ-переходу пропорційна квадрату норми похідної дипольного моменту:
    \begin{equation}
        I_k = 974.8638 \times \left|\frac{\partial \boldsymbol{\mu}}{\partial Q_k}\right|^2_{\text{a.u.}}
        \quad \text{[км/моль]}
    \end{equation}

    \begin{minted}{python}
    intensity = np.linalg.norm(dip_derivative)**2 * 974.8638
    intensities.append(intensity)
    \end{minted}
\end{enumerate}



Хоча метод простий, він вимагає виконання двох \texttt{SCF}-розрахунків на кожну моду
і може бути чутливим до вибору~$\delta$.

%---------------------------------------------------------
\subsection{Аналітичний підхід у \texttt{PySCF}}
%---------------------------------------------------------
Модуль \inlinecode{pyscf.prop.freq} реалізує аналітичний обрахунок похідних
дипольного моменту по атомних координатах $\partial\boldsymbol{\mu}/\partial \mathbf{R}$.

\paragraph{Перехід до нормальних координат.}
Похідна дипольного моменту за нормальною координатою $Q_i$ виражається через
матрицю нормальних мод $\mathbf{L}$:
\[
\frac{\partial \vect{\mu}}{\partial Q_i}
= \sum_{A=1}^{N_{\text{атм}}}
    \frac{\partial \vect{\mu}}{\partial \vect{R}_{A}} \cdot \frac{\partial \vect{R}_A}{\partial Q_i}
= \sum_{A=1}^{N_{\text{атм}}}
    \frac{\partial \vect{\mu}}{\partial \vect{R}_{A}} \cdot \vect{L}_A^i,
\]
де $\vect L_A^i = (L_{Ax}^i, L_{Ay}^i, L_{Az}^i)$ --- вектор нормальної моди для атома $A$.

\paragraph{Структура тензора похідних.}
У \texttt{PySCF} тензор похідних \inlinecode{vib.de} має розмірність $(N_{\text{атм}}, N_{\text{атм}}, 3, 3)$
і зберігає похідні компонент дипольного моменту за координатами атомів:
\[
(\texttt{vib.de})_{ABaj} = \frac{\partial \mu_j^{(A)}}{\partial R_{Ba}},
\]

Діагональні елементи ($A=B$) відповідають похідній внеску атома в дипольний момент по його власних координатах:
\[
D_{Aaj} = (\texttt{vib.de})_{AAaj}
= \frac{\partial \mu_j}{\partial R_{Aa}}.
\]

\begin{minted}{python}
# Діагональний тензор: (N_atm, 3, 3)
de_diag = np.einsum('AAaj->Aaj', vib.de)
\end{minted}

Для кожного атома $A$ отримуємо матрицю Якобі дипольного моменту:
\[
\mathbf{D}_A =
\begin{pmatrix}
\partial \mu_x/\partial x_A & \partial \mu_y/\partial x_A & \partial \mu_z/\partial x_A \\
\partial \mu_x/\partial y_A & \partial \mu_y/\partial y_A & \partial \mu_z/\partial y_A \\
\partial \mu_x/\partial z_A & \partial \mu_y/\partial z_A & \partial \mu_z/\partial z_A
\end{pmatrix}.
\]

\paragraph{Похідні за нормальними координатами.}
Матриця похідних дипольного моменту за всіма нормальними модами:
\[
\frac{\partial \mu_j}{\partial Q_i}
= \sum_{A=1}^{N_{\text{атм}}} \sum_{a=x,y,z} L_{Aa}^i \, D_{Aaj}.
\]

\begin{minted}{python}
# Тензорна згортка: (N_modes, 3)
dipole_derivs = np.einsum('iAa,Aaj->ij', modes, de_diag)
\end{minted}

\paragraph{Інтенсивності ІЧ-переходів.}
Інтенсивність $i$-го коливального переходу:
\[
I_i = 974.8638 \sum_{j=x,y,z} \left( \frac{\partial \mu_j}{\partial Q_i} \right)^2
= 974.8638 \, \left\| \frac{\partial \boldsymbol{\mu}}{\partial Q_i} \right\|^2
\]

\begin{minted}{python}
# Інтенсивності для всіх мод
ir_intensities = 974.8638 * np.sum(dipole_derivs**2, axis=1)
\end{minted}

І весь код разом:

\begin{minted}{python}
import numpy as np
from pyscf import gto, scf
from pyscf.prop import freq

mol = gto.M(atom='O 0 0 0; H 0 -0.757 0.587; H 0 0.757 0.587', basis='6-31g')
mf = scf.RHF(mol).run()

vib = freq.rhf.Frequency(mf)
freqs, modes = vib.kernel()
natm = mol.natm

# Діагональні елементи
de_diag = np.einsum('AAaj->Aaj', vib.de)  # (natm, 3, 3)

# Всі дипольні похідні одразу
dipole_derivs = np.einsum('iAa,Aaj->ij', modes, de_diag)

# Всі інтенсивності одразу
ir_intensities = 974.8638 * np.sum(dipole_derivs**2, axis=1)

# Вивід
print("\n" + "="*70)
print("  ІЧ СПЕКТР H₂O (PySCF, оптимізовано)")
print("="*70)
print(f"{'Мода':<8}{'Частота (см⁻¹)':>18}{'Інтенсивність (km/mol)':>28}")
print("-"*70)
for i in range(6, len(freqs)):
    print(f"{i-5:<8d}{freqs[i]:>18.2f}{ir_intensities[i]:>28.2f}")
print("="*70)
\end{minted}

На відміну від чисельного підходу, цей метод:
\begin{itemize}
    \item не потребує перерахунку енергії при зсуві атомів;
    \item базується на аналітичних градієнтах і похідних;
    \item забезпечує вищу точність і швидкість.
\end{itemize}


\subsubsection{Результати для \ce{H2O}}

\paragraph{Нормальні моди води.}

Експериментальні дані

\begin{table}[h]
\centering
\small
\caption{Експериментальні дані коливань молекули H$_2$O (джерела: \url{https://cccbdb.nist.gov/expvibs2x.asp}).}
\label{tab:exp_vib_H2O}
\begin{tblr}{
colspec = {Q[c,m] Q[l,m,] Q[c,m] Q[c,m] Q[c,m] X[l,m]},
row{1} = {font=\bfseries, c},
hline{1,2,Z} = {-}{0.08em},
}
Мода & Симетрія & {Частота\\ гармонічна (см$^{-1}$)} & {Інтенсивність\\ (км моль$^{-1}$)} \\
1 & \ce{H-O-H} вигин &  1648.5 & 62.5  \\
2 & \ce{O-H} сим. розтяг &  3832.2 & 2.93 \\
3 & \ce{O-H} антисим. розтяг  & 3942.5 & 41.7  \\
\end{tblr}
\end{table}


Молекула H$_2$O має $3N - 6 = 3$ коливальні моди:

\begin{center}
\begin{tblr}{
colspec = {Q[c,m] Q[c,m] Q[c,m] Q[c,m]},
row{1} = {font=\bfseries, c, m},
hline{1, Z} = {solid},
}
Мода & $\nu$ (см$^{-1}$) & $I$ (км/моль) & Тип \\
1 & 1828.9 & 1.4 & Слабка (деформаційна) \\
2 & 3906.0 & 102.3 & Сильна (симетричне валентне) \\
3 & 4001.1 & 138.6 & Сильна (асиметричне валентне) \\
\end{tblr}
\end{center}

\smallskip
\textit{Примітка.} Усі три моди молекули \ce{H2O} є ІЧ-активними
(група симетрії $C_{2v}$). Найсильніше поглинання відповідає
валентним коливанням, що зумовлюють інтенсивну ІЧ-смугу пари води
в атмосфері (важливий внесок у парниковий ефект).

\subsection*{Резюме}

У цьому розділі розглянуто теоретичні основи та практичні методи розрахунку
інфрачервоних спектрів молекул.

\paragraph{Ключові положення:}
\begin{itemize}
    \item ІЧ-інтенсивність коливальної моди визначається квадратом похідної
    дипольного моменту за нормальною координатою:
    $I_i = 974.8638 \left|\frac{\partial\boldsymbol{\mu}}{\partial Q_i}\right|^2$ [км/моль].

    \item Інтенсивність відображає \emph{силу взаємодії} молекули з ІЧ-випромінюванням:
    чим більша зміна дипольного моменту при коливанні, тим інтенсивніша смуга поглинання.
    Класифікація: сильні смуги ($I > 100$~км/моль), середні ($10 < I \leq 100$),
    слабкі ($I \leq 10$~км/моль).

    \item Два методи обчислення похідних:
    \begin{itemize}
        \item \textit{Чисельний} --- простий, але вимагає додаткових SCF-розрахунків
        для зміщених геометрій та чутливий до вибору кроку $h$;
        \item \textit{Аналітичний} --- ефективніший і точніший, базується на аналітичних
        градієнтах густини заряду. Реалізований у \texttt{PySCF} як для HF,
        так і для DFT-методів.
    \end{itemize}

    \item Для молекули H$_2$O розраховані частоти ($\sim 1829$, $3906$, $4001$~см$^{-1}$)
    якісно узгоджуються з експериментом ($1648$, $3832$, $3943$~см$^{-1}$),
    а систематичне завищення пояснюється гармонічним наближенням і вибором базису.

    \item Інтенсивності ($I_1 = 1.4$, $I_2 = 102.3$, $I_3 = 138.6$~км/моль)
    демонструють, що найінтенсивніші смуги відповідають валентним коливанням \ce{O-H} зв'язків,
    тоді як деформаційна мода має слабку інтенсивність. Експериментальні значення
    ($62.5$, $2.93$, $41.7$~км/моль) показують іншу картину через ангармонічні ефекти
    та обмеження гармонічного наближення.

    \item Висока інтенсивність валентних коливань робить пару води потужним поглиначем
    ІЧ-випромінювання в атмосфері, що має значення для клімату та парникового ефекту.
\end{itemize}

\paragraph{Практичне значення:}
Розрахунки ІЧ-спектрів є важливим інструментом для:
\begin{itemize}
    \item ідентифікації функціональних груп та структури молекул;
    \item інтерпретації експериментальних спектрів та передбачення відносних інтенсивностей;
    \item передбачення спектральних характеристик нових сполук;
    \item дослідження міжмолекулярних взаємодій та динаміки;
    \item кількісного аналізу концентрацій речовин за допомогою закону Бугера-Ламберта-Бера.
\end{itemize}


